\chapter{The MSCKF Visual Update}
\label{ch:msckfupdate}

\begin{objectives}
\begin{itemize}
  \item Build the linearized measurement model of a feature track over the clones that observed it.
  \item Eliminate the feature by left-nullspace projection, and prove that no information about the poses is lost.
  \item Compute the projection efficiently with Householder reflections.
  \item Gate features, stack them, compress the stack with QR, and apply the update to the IMU state and all clones.
\end{itemize}
\end{objectives}

% ------------------------------------------------------------------------------
\section{A feature track}

A feature $j$ has been observed at the clones $\mathcal{S}_j = \{i_1, \dots, i_M\}$, with undistorted normalized observations
$z_{j,i}$. The update is triggered when the track is lost or when its oldest observation's clone is about to be marginalized,
and requires $M \ge 3$ (AGENTS.md F5). The first step is to estimate the feature position $\hat p_f$ in the world frame by
triangulation from the clone camera poses
\begin{equation}
  \hat R_{wc_i} = \hat R_{wb_i}R_{bc}, \qquad \hat p_{wc_i} = \hat p_{wb_i} + \hat R_{wb_i}p_{bc},
\end{equation}
using the linear initialization and Gauss--Newton refinement of \cref{sec:triangulation}, with the positive-depth, parallax and
reprojection checks. Features that fail are discarded.

% ------------------------------------------------------------------------------
\section{Residuals and Jacobians}

For each observation, the predicted measurement uses the clone estimate and the triangulated point,
\begin{equation}
  \hat p_{c_i} = R_{bc}\T\big(\hat R_{wb_i}\T(\hat p_f - \hat p_{wb_i}) - p_{bc}\big), \qquad
  r_{j,i} = z_{j,i} - \pi(\hat p_{c_i}) .
\end{equation}
Linearizing with respect to the error state $\delta x$ (IMU and clones) and the feature position error $\delta p_f = p_f - \hat p_f$,
\begin{equation}
  r_{j,i} \approx H_{x,ji}\,\delta x + H_{f,ji}\,\delta p_f + n_{j,i}, \qquad n_{j,i} \sim \N(0, \sigma_{uv}^2I_2).
\end{equation}
The structure of these Jacobians follows from the chain rule
\begin{equation}
  H_{x,ji} = \frac{\partial\pi}{\partial p_c}\Big|_{\hat p_{c_i}}\frac{\partial p_{c_i}}{\partial\delta x},\qquad
  H_{f,ji} = \frac{\partial\pi}{\partial p_c}\Big|_{\hat p_{c_i}}\frac{\partial p_{c_i}}{\partial p_f},
\end{equation}
with the projection Jacobian of \cref{prop:proj-jacobian}. The observation at clone $i$ depends only on that clone's pose, so
$H_{x,ji}$ is zero except in the six columns of $(\delta p_i, \delta\theta_i)$. In particular it is zero in all IMU columns.

\begin{derivex}{the visual measurement Jacobians $H_{x,ji}$ and $H_{f,ji}$}{sol:Hcam}
With $p_{c_i} = R_{bc}\T\big(R_{wb_i}\T(p_f - p_{wb_i}) - p_{bc}\big)$, the local perturbation $R_{wb_i} = \hat R_{wb_i}\Exp(\delta\theta_i)$ and
$p_{wb_i} = \hat p_{wb_i} + \delta p_i$:
\begin{enumerate}
  \item Derive $\partial p_{c_i}/\partial\delta\theta_i$ with the inverse-action row of \cref{tab:jacobians}. Be careful about which
        vector the rotation acts on.
  \item Derive $\partial p_{c_i}/\partial\delta p_i$ and $\partial p_{c_i}/\partial p_f$. How are they related, and why does that make physical sense?
  \item Multiply by $\partial\pi/\partial p_c$ and place the blocks in the frozen column layout. State the sparsity pattern of the stacked
        $H_{x,j}$.
\end{enumerate}
\textbf{Check:} Milestone 10b requires finite-difference tests of the projection Jacobians with respect to the clone
$\delta p$, the clone $\delta\theta$ and $p_f$, at a non-trivial clone pose with non-trivial extrinsics.
\end{derivex}

Stacking the $M$ observations of feature $j$,
\begin{equation}
  r_j \approx H_{x,j}\,\delta x + H_{f,j}\,\delta p_f + n_j, \qquad
  r_j, n_j \in \R^{2M},\quad H_{x,j} \in \R^{2M\times(15+6N)},\quad H_{f,j} \in \R^{2M\times3}.
  \label{eq:stacked-feature}
\end{equation}

\begin{figure}[t]
\centering
\begin{tikzpicture}[scale=0.42]
  % H_x
  \draw[black!40] (0,0) rectangle (15,-10);
  \fill[black!8] (0,0) rectangle (3.75,-10);
  \node[font=\scriptsize] at (1.875,0.6) {IMU (15)};
  \foreach \k in {0,...,4} {
    \fill[primary!70] ({3.75 + 1.5*(\k+2)}, {-2*\k}) rectangle ({3.75 + 1.5*(\k+3)}, {-2*\k-2});
  }
  \foreach \k in {0,...,6} \draw[black!25] ({3.75 + 1.5*\k}, 0) -- ({3.75 + 1.5*\k}, -10);
  \node[font=\scriptsize] at (9.4,0.6) {clones (6 columns each)};
  \node[font=\small] at (7.5,-11) {$H_{x,j}$};
  % H_f
  \draw[black!40] (16.5,0) rectangle (18,-10);
  \fill[plotgreen!70] (16.5,0) rectangle (18,-10);
  \node[font=\small] at (17.25,-11) {$H_{f,j}$};
  % r
  \draw[black!40] (19.5,0) rectangle (20.2,-10);
  \fill[plotorange!60] (19.5,0) rectangle (20.2,-10);
  \node[font=\small] at (19.85,-11) {$r_j$};
  \draw[decorate, decoration={brace, amplitude=4pt}] (-0.4,-10) -- (-0.4,0) node[midway, left=4pt, font=\scriptsize] {$2M$ rows};
\end{tikzpicture}
\caption{Sparsity of the stacked feature Jacobian for a feature seen by five consecutive clones. Each pair of rows depends
on one clone; the IMU columns are zero. $H_{f,j}$ is dense.}
\label{fig:hx-sparsity}
\end{figure}

\section{The problem with the feature}

\Cref{eq:stacked-feature} cannot be used directly in an update. The feature error $\delta p_f$ is not part of the state, so there is
nothing to correct. Worse, $\hat p_f$ was computed from the same measurements $z_{j,i}$ and the same clone estimates, so its error is
correlated with both $n_j$ and $\delta x$ in a way the filter's covariance does not describe. Treating $\hat p_f$ as known would
double count the measurements. EKF-SLAM resolves this by putting $p_f$ in the state. The MSCKF eliminates it.

% ------------------------------------------------------------------------------
\section{Left-nullspace projection}

$H_{f,j}$ is a tall $2M\times3$ matrix. If the feature is well triangulated it has full column rank 3, and its left nullspace
has dimension $2M - 3$ (\cref{sec:subspaces}). Let $A_j \in \R^{2M\times(2M-3)}$ have orthonormal columns spanning
$\nullsp(H_{f,j}\T)$, so that $A_j\T H_{f,j} = 0$ and $A_j\T A_j = I$. Multiplying \eqref{eq:stacked-feature} by $A_j\T$,
\begin{equation}
  r_{o,j} := A_j\T r_j \approx \underbrace{A_j\T H_{x,j}}_{H_{o,j}}\,\delta x + \underbrace{A_j\T H_{f,j}}_{=\,0}\,\delta p_f + A_j\T n_j
  = H_{o,j}\,\delta x + n_{o,j} .
  \label{eq:projected}
\end{equation}
The feature has disappeared, and the projected noise is still white with the same variance:
$\Cov(n_{o,j}) = \sigma_{uv}^2A_j\T A_j = \sigma_{uv}^2 I_{2M-3}$. \Cref{eq:projected} is a standard EKF measurement of the state alone,
of dimension $2M - 3$.

\begin{figure}[t]
\centering
\begin{tikzpicture}[scale=1.0]
  \fill[plotgreen!15] (-2.8,-0.9) -- (1.8,-0.9) -- (3.0,0.9) -- (-1.6,0.9) -- cycle;
  \node[plotgreen!60!black, font=\footnotesize, align=center] at (-1.6,-0.55) {$\range(H_f)$: residual directions\\ explainable by moving the feature};
  \draw[->, thick, primary] (0,0) -- (0,2.6) node[above, font=\footnotesize] {$\nullsp(H_f\T)$: what only the poses can explain};
  \draw[->, very thick, plotorange] (0,0) -- (1.6,1.9) node[right, font=\footnotesize] {$r_j$};
  \draw[dashed, black!50] (1.6,1.9) -- (0,1.9);
  \draw[->, very thick, conv] (0,0) -- (0,1.9) node[left, font=\footnotesize] {$A_jA_j\T r_j$};
  \draw[dashed, black!50] (1.6,1.9) -- (1.6,0.0);
  \fill (0,0) circle (1.5pt);
\end{tikzpicture}
\caption{Geometry of the nullspace projection. Residual components inside the column space of $H_f$ can be absorbed by an
adjustment of the unknown feature position, so they say nothing about the poses. The component orthogonal to it, in the
left nullspace, cannot be explained by any feature position and must be explained by the poses.}
\label{fig:nullspace}
\end{figure}

\begin{theorem}[Nullspace projection loses no information]
\label{thm:nullspace-optimal}
Suppose $\delta p_f$ has no prior information and $H_{f,j}$ has full column rank. Then the information about $\delta x$ in the
projected measurement \eqref{eq:projected} equals the information about $\delta x$ in \eqref{eq:stacked-feature} after $\delta p_f$ is marginalized
out:
\begin{equation}
  H_{o,j}\T H_{o,j} = H_{x,j}\T\Big(I - H_{f,j}(H_{f,j}\T H_{f,j})\inv H_{f,j}\T\Big)H_{x,j},
\end{equation}
and likewise for the information vectors $H_{o,j}\T r_{o,j}$. The two formulations therefore produce identical updates.
\end{theorem}
\begin{proof}
The joint information matrix of $(\delta x, \delta p_f)$ contributed by \eqref{eq:stacked-feature} is
$\sigma^{-2}\begin{bmatrix} H_x\T H_x & H_x\T H_f \\ H_f\T H_x & H_f\T H_f\end{bmatrix}$. Marginalizing $\delta p_f$ in information form is the Schur
complement (\cref{thm:gauss-cond}), $\sigma^{-2}\big(H_x\T H_x - H_x\T H_f(H_f\T H_f)\inv H_f\T H_x\big)$, and by the projector identity
\eqref{eq:nullspace-projector}, $I - H_f(H_f\T H_f)\inv H_f\T = A A\T$. The same algebra applies to the information vector
$\sigma^{-2}H_x\T(I - H_f(H_f\T H_f)\inv H_f\T)r$.
\end{proof}

An equivalent way to see this: complete $A_j$ to an orthogonal matrix $Q = [Q_1\ \ A_j]$. Multiplying by $Q\T$ splits the equations
into three that involve $\delta p_f$ (through the invertible $3\times3$ matrix $Q_1\T H_f$) and $2M - 3$ that do not. Since $\delta p_f$ is free,
the first three can always be satisfied by choosing $\delta p_f$, so they carry no information about $\delta x$. The rest carry all
of it.

\begin{remark}[how many observations]
Two observations ($M = 2$) leave $2M - 3 = 1$ constraint, essentially the epipolar constraint between two clones. It is weak and
sensitive to triangulation error, so the project requires $M \ge 3$ (AGENTS.md F5), which gives at least three constraints and
a better-conditioned triangulation. If the triangulation is poor (low parallax), $H_{f,j}$ is close to rank-deficient, the
projection is ill-defined, and the feature should have been rejected by the triangulation checks.
\end{remark}

\subsection{Computing the projection}

$A_j$ is never formed. By \cref{prop:qr-nullspace}, the last $2M - 3$ columns of the $Q$ factor of a full QR decomposition of $H_{f,j}$
span its left nullspace. Householder QR of the $2M\times3$ matrix $H_{f,j}$ uses exactly three reflections. Applying the same three
reflections to $H_{x,j}$ and $r_j$, and keeping the last $2M - 3$ rows of the results, gives $H_{o,j}$ and $r_{o,j}$ directly. The cost is
$\mathcal{O}(M\cdot k)$ for $k$ columns, and only the $6M$ non-zero clone columns of $H_{x,j}$ need processing. Givens rotations achieve the
same with finer control over sparsity and were suggested in the original paper. Milestone 10b checks
$\|A\T H_f\| < 10^{-10}$ and the orthonormality of $A$ in a test, where forming $A$ explicitly is acceptable.

% ------------------------------------------------------------------------------
\section{Gating, stacking and compression}

\paragraph{Per-feature gating.} Each projected residual is gated before it is accepted (AGENTS.md F5):
\begin{equation}
  \gamma_j = r_{o,j}\T\big(H_{o,j}PH_{o,j}\T + \sigma_{uv}^2I\big)\inv r_{o,j} < \chi^2_{0.95}(2M - 3).
\end{equation}
This rejects tracks that passed the front end's checks but are inconsistent with the filter's current belief: tracks on moving
objects, slowly drifting tracks, and features with wrong associations.

\paragraph{Stacking.} The accepted features are stacked into one measurement,
\begin{equation}
  r_o = H_o\,\delta x + n_o, \qquad r_o \in \R^d,\quad d = \sum_j (2M_j - 3),
\end{equation}
with the number of rows capped at a compile-time maximum so that the stacked workspace is fixed-size. With 50 features of 6
observations each, $d = 450$, much larger than the state dimension $n = 75$.

\paragraph{QR compression.} When $d > n$, the update can be made cheaper without approximation.

\begin{proposition}[Measurement compression]
\label{prop:compression}
Let $H_o = Q_1T_H$ be a thin QR decomposition, with $T_H \in \R^{n\times n}$ upper triangular, and define $r_n = Q_1\T r_o$. Then the
measurement $r_n = T_H\,\delta x + n_n$ with $n_n \sim \N(0, \sigma_{uv}^2I_n)$ yields exactly the same update as $(r_o, H_o)$.
\end{proposition}
\begin{proof}
The update depends on the measurement only through the information $H_o\T H_o/\sigma^2$ and the information vector $H_o\T r_o/\sigma^2$
(\cref{eq:kf-info}). Since $Q_1\T Q_1 = I$, $T_H\T T_H = H_o\T Q_1Q_1\T H_o = H_o\T H_o$ (as $Q_1Q_1\T$ projects onto $\range(H_o)$, which contains
the columns of $H_o$), and $T_H\T r_n = H_o\T Q_1Q_1\T r_o = H_o\T r_o$. The discarded components $Q_2\T r_o$ are pure noise, independent of $\delta x$.
\end{proof}

Compression reduces the innovation covariance from $d\times d$ to at most $n\times n$, and the Joseph update from
$\mathcal{O}(dn^2 + d^3)$ to $\mathcal{O}(n^3)$ after an $\mathcal{O}(dn^2)$ QR. Only the non-zero columns of $H_o$ (the clone columns) need to
take part in the decomposition. Compression is optional in the project (AGENTS.md F5); it is worthwhile when $d \gg n$.

% ------------------------------------------------------------------------------
\section{The update}

The stacked (and possibly compressed) measurement is applied with the generic update of \cref{ch:updates}: $S = HPH\T +
\sigma_{uv}^2I$, gain by Cholesky solve, Joseph form, then injection into the IMU state and \emph{every} clone, normalization of every
touched quaternion, and the reset with one $I - \tfrac12\sx{\delta\hat\theta_i}$ block per attitude. Corrections to the clones are
real: they improve the poses that future features will be triangulated from.

\begin{algorithm}[t]
\caption{MSCKF visual update}
\label{alg:msckf-update}
\begin{algorithmic}[1]
\Require filter state with clones, set of feature tracks ready for update
\State $d \gets 0$
\For{each ready feature $j$ with $M_j \ge 3$ observations}
  \State triangulate $\hat p_f$ (linear + Gauss--Newton); \textbf{if} checks fail \textbf{then} discard, \textbf{continue}
  \State build $r_j$, $H_{x,j}$ (clone columns only), $H_{f,j}$
  \State apply Householder QR of $H_{f,j}$ to $[H_{x,j}\ \ r_j]$; keep rows $4\dots2M_j$ $\to$ $H_{o,j}$, $r_{o,j}$
  \If{$\gamma_j < \chi^2_{0.95}(2M_j - 3)$ \textbf{and} $d + 2M_j - 3 \le d_{max}$}
    \State append $H_{o,j}$, $r_{o,j}$ to the stacked workspace; $d \gets d + 2M_j - 3$
  \Else \State count the rejection \EndIf
  \State mark the track's observations as consumed
\EndFor
\If{$d > 0$}
  \State optionally compress with QR if $d > n$
  \State EKF update (Joseph), inject into IMU and all clones, normalize, reset with $G$
\EndIf
\end{algorithmic}
\end{algorithm}

\begin{projectlink}[Milestone 10b]
The exit criteria test every link of \cref{alg:msckf-update}: projection Jacobians by finite differences; triangulation exact on
noise-free data and bounded under noise, rejecting negative-depth and low-parallax features; $\|A\T H_f\| < 10^{-10}$ with
orthonormal $A$; a full synthetic VIO run without divergence, with NEES reported (yaw inconsistency expected and documented);
zero allocations with fixed-max workspaces; and MH\_01 end to end with ATE reported.
\end{projectlink}

\begin{takeaways}
\begin{itemize}
  \item Each observation depends on one clone; $H_{x,j}$ is zero in the IMU columns and dense in the feature's $H_{f,j}$.
  \item Projecting onto the left nullspace of $H_{f,j}$ removes the feature, keeps the noise white, and loses no information about
        the poses (it equals Schur-complement marginalization).
  \item Three Householder reflections applied to $[H_{x,j}\ r_j]$ compute the projection without forming $A_j$.
  \item Gate each feature with $\chi^2_{0.95}(2M - 3)$, stack, compress with QR if $d > n$, and update IMU and clones together.
\end{itemize}
\end{takeaways}

\begin{exercises}
  \item Verify \cref{thm:nullspace-optimal} numerically: compare the posterior covariance of $\delta x$ from the projected update with
        that from an update of the augmented state $(\delta x, \delta p_f)$ with a very large prior on $\delta p_f$.
  \item Show that the three discarded rows $Q_1\T r_j$ determine the correction to $\hat p_f$ given $\delta x$. How would you recover an
        updated feature estimate if you wanted one?
  \item For $M = 2$, compute $A_j$ explicitly and interpret the single projected residual geometrically.
  \item Implement the Householder application of \cref{alg:msckf-update} and compare its cost and result with forming $A_j$ from a full QR.
  \item Prove that the projected noise vectors $n_{o,j}$ of two different features are independent if their raw noises are,
        so that stacking them with a block-diagonal (here scalar) noise covariance is correct. Why are the projected
        \emph{residuals} themselves correlated?
\end{exercises}
